% Embedded Systems lab report -- compile with pdflatex (twice, for the ToC)
\documentclass[11pt,a4paper]{article}

\usepackage[margin=2.5cm]{geometry}
\usepackage{lmodern}      % scalable Type1 fonts for T1; without this,
\usepackage[T1]{fontenc}  % MiKTeX builds slow, fuzzy METAFONT bitmaps
\usepackage[utf8]{inputenc}
\usepackage{graphicx}
\usepackage{booktabs}
\usepackage{xcolor}
\usepackage{listings}
\usepackage[hidelinks]{hyperref}

% --- C code listing style -------------------------------------------------
\definecolor{codegray}{gray}{0.96}
\definecolor{codegreen}{rgb}{0.0,0.5,0.0}
\definecolor{codeblue}{rgb}{0.0,0.0,0.6}

\lstdefinestyle{cstyle}{
  language=C,
  backgroundcolor=\color{codegray},
  basicstyle=\ttfamily\footnotesize,
  keywordstyle=\color{codeblue}\bfseries,
  commentstyle=\color{codegreen}\itshape,
  stringstyle=\color{red!70!black},
  numbers=left,
  numberstyle=\tiny\color{gray},
  numbersep=8pt,
  breaklines=true,
  frame=single,
  framerule=0pt,
  showstringspaces=false,
  tabsize=2,
  captionpos=b,
}
\lstset{style=cstyle}

% --- Title ----------------------------------------------------------------
% >>> FILL THIS IN <<<
\newcommand{\studentid}{R12942157}

\title{Embedded Systems Laboratory Report\\
       \large FreeRTOS Task Synchronisation on the B-L475E-IOT01A\\
       \normalsize Basic Problem and Option Problem A}
\author{YU-CHEN, CHU\\
        \normalsize Student ID: \studentid}
\date{\today\\[2ex]
      \normalsize Source code:
      \url{https://github.com/oldpig123/embedded_system/tree/master/lab1}}

\begin{document}
\maketitle
\tableofcontents
\newpage

% =========================================================================
\section{Introduction}

The purpose of this exercise was to implement two independent blinking
behaviours on a single LED, each driven by its own FreeRTOS task, and to
coordinate them so that the two behaviours never interfere with one another.
Because both tasks manipulate the same physical output pin, the exercise is
fundamentally about \emph{synchronisation}: deciding which task is allowed to
drive the LED at any given moment, and how an asynchronous external event
(a button press, or a periodic timer expiry) is communicated to the task that
should react to it.

Two related but distinct mechanisms are used throughout. A \emph{binary
semaphore} and a \emph{message queue} are used for \emph{signalling}: telling a
sleeping task that an event has occurred. A \emph{mutex} is used for
\emph{mutual exclusion}: guaranteeing that only one task drives the shared LED
at a time. Distinguishing these two roles turned out to be the central
conceptual lesson of the exercise.

The extension (Option Problem A) adds the requirement that the button
distinguish a normal press from a long press, which forces the signalling
mechanism to carry \emph{data} rather than a bare notification, and therefore
replaces the semaphore with a message queue.

% =========================================================================
\section{Development Environment}

\begin{table}[h]
\centering
\begin{tabular}{ll}
\toprule
Item & Detail \\
\midrule
Board            & B-L475E-IOT01A Discovery kit \\
MCU              & STM32L475VG (ARM Cortex-M4F) \\
RTOS             & FreeRTOS, accessed through the CMSIS-RTOS v2 API \\
Code generation  & STM32CubeMX \\
Toolchain        & \texttt{arm-none-eabi-gcc}, CMake build \\
Indicator        & LED2 on PB14 \\
Input            & User button on PC13, via EXTI line 13 \\
\bottomrule
\end{tabular}
\caption{Development environment.}
\end{table}

All RTOS objects are created with the CMSIS-RTOS v2 API
(\texttt{osThreadNew}, \texttt{osSemaphoreNew}, \texttt{osMutexNew},
\texttt{osTimerNew}, \texttt{osMessageQueueNew}) rather than calling the
underlying FreeRTOS functions directly.

% =========================================================================
\section{Basic Problem}

\subsection{Requirements}

Two tasks share LED2 and must not interfere with each other:

\begin{itemize}
  \item \textbf{Task\_1} --- pressing the user button triggers LED2 to blink for
        5\,s at 1\,Hz. The EXTI interrupt handler detects the press and informs
        Task\_1 through a binary semaphore.
  \item \textbf{Task\_2} --- LED2 blinks for 2\,s at 10\,Hz, repeating every
        10\,s. A timer interrupt detects the time-up event and informs Task\_2
        through a second binary semaphore.
\end{itemize}

A mutex or semaphore must protect the blinking procedures so that while one
blinking sequence is in progress, the other waits.

\subsection{System Architecture}

Both tasks follow the same three-stage pattern:

\begin{center}
\texttt{external event} $\rightarrow$ \texttt{ISR / callback} $\rightarrow$
\texttt{signalling object} $\rightarrow$ \texttt{task performs blinking}
\end{center}

The interrupt or callback performs no work beyond signalling; all timing and
GPIO manipulation happens in task context. This separation is essential,
because delaying inside an interrupt handler is either unsafe or ineffective
(Section~\ref{sec:problems}).

\begin{table}[h]
\centering
\begin{tabular}{lllll}
\toprule
Role & Function & Priority & Trigger & Blink pattern \\
\midrule
Task\_1 & \texttt{StartTask03}      & \texttt{osPriorityHigh7} & EXTI13 ISR        & 5\,s @ 1\,Hz \\
Task\_2 & \texttt{StartDefaultTask} & \texttt{osPriorityNormal}& periodic RTOS timer & 2\,s @ 10\,Hz \\
\bottomrule
\end{tabular}
\caption{Task structure. A third thread, \texttt{StartTask02}, remains from an
earlier exercise and is idle.}
\end{table}

\subsection{Synchronisation Objects}

\begin{table}[h]
\centering
\begin{tabular}{lll}
\toprule
Object & Creation & Purpose \\
\midrule
\texttt{myMutex01Handle}      & \texttt{osMutexNew(NULL)}                  & protects LED2 \\
\texttt{myBinarySem02Handle}  & \texttt{osSemaphoreNew(1, 0, NULL)}        & timer $\rightarrow$ Task\_2 \\
\texttt{myTimer01Handle}      & \texttt{osTimerNew(\ldots, osTimerPeriodic, \ldots)} & 10\,s periodic tick \\
\texttt{myQueue01Handle}      & \texttt{osMessageQueueNew(16, sizeof(uint32\_t), NULL)} & ISR $\rightarrow$ Task\_1 \\
\bottomrule
\end{tabular}
\caption{RTOS objects. The queue replaces Task\_1's binary semaphore in
Option Problem A.}
\end{table}

All objects are created in \texttt{main()} after \texttt{osKernelInitialize()}
and before \texttt{osKernelStart()}. Both semaphores are created with an
initial count of zero, so that the waiting task genuinely blocks until a real
event occurs rather than passing through immediately on its first attempt.

\subsection{Timing Design}

A blink of frequency $f$ has period $T = 1/f$, and a 50\,\% duty cycle is
obtained by toggling the pin every $T/2$. The number of loop iterations
required for a total duration $D$ is therefore $N = D/T$.

\begin{table}[h]
\centering
\begin{tabular}{lllll}
\toprule
Sequence & Frequency & Half period & Duration & Iterations \\
\midrule
Task\_2 burst          & 10\,Hz & 50\,ms  & 2\,s & 20 \\
Task\_1 normal press   & 1\,Hz  & 500\,ms & 5\,s & 5 \\
Task\_1 long press     & 10\,Hz & 50\,ms  & 5\,s & 50 \\
\bottomrule
\end{tabular}
\caption{Blink timing parameters.}
\end{table}

Each loop iteration contains two toggles and two delays, so one iteration
corresponds to one full on/off cycle.

\subsection{Implementation}

The timer callback does nothing but release the semaphore that Task\_2 is
waiting on:

\begin{lstlisting}[caption={Periodic timer callback signalling Task\_2.}]
void Timer01_Callback(void *argument)
{
  osSemaphoreRelease(myBinarySem02Handle);
}
\end{lstlisting}

Task\_2 blocks on that semaphore and performs its burst when signalled. Note
that the mutex is acquired and released around each individual toggle rather
than around the whole burst; the reason is discussed in
Section~\ref{sec:granularity}.

\begin{lstlisting}[caption={Task\_2: 2\,s at 10\,Hz, triggered by the timer.}]
void StartDefaultTask(void *argument)
{
  for(;;)
  {
    osSemaphoreAcquire(myBinarySem02Handle, osWaitForever);
    for(int i = 0; i < 20; i++)
    {
      osMutexAcquire(myMutex01Handle, osWaitForever);
      HAL_GPIO_TogglePin(GPIOB, LED2_Pin);
      osMutexRelease(myMutex01Handle);
      osDelay(50);
      osMutexAcquire(myMutex01Handle, osWaitForever);
      HAL_GPIO_TogglePin(GPIOB, LED2_Pin);
      osMutexRelease(myMutex01Handle);
      osDelay(50);
    }
  }
}
\end{lstlisting}

Because the task no longer performs its own inter-burst delay, the 10-second
period is governed entirely by the timer. The task simply returns to blocking
on the semaphore as soon as a burst finishes.

\subsection{Results}

The observed behaviour is a 2-second burst of fast blinking followed by
8 seconds of darkness, repeating indefinitely --- a 10-second period, as
required.

One initially surprising detail is that nothing happens for the first
10 seconds after reset. This is correct: the periodic timer starts counting
when \texttt{osTimerStart()} executes, so its first expiry occurs a full period
later. Until then Task\_2 is blocked on a semaphore whose count is zero.
Should an immediate first burst be desired, releasing the semaphore once
manually before \texttt{osKernelStart()} would shift the phase without altering
the period.

% =========================================================================
\section{Option Problem A: Long Press Detection}

\subsection{Requirements}

The button must distinguish two events: a normal press, and a long press in
which the button is held for more than one second. A normal press produces the
original 5\,s at 1\,Hz blink; a long press produces a 5\,s at 10\,Hz blink.
Because Task\_1 must now be told \emph{which} of two events occurred, the
specification requires a message queue rather than a binary semaphore.

\subsection{Detecting Press Duration}

Measuring how long a button is held requires knowing both when it was pressed
and when it was released. The original configuration generated an interrupt
only on the falling edge, so the release was invisible to software. The EXTI
configuration for PC13 was therefore changed from
\texttt{GPIO\_MODE\_IT\_FALLING} to \texttt{GPIO\_MODE\_IT\_RISING\_FALLING}
in STM32CubeMX and the project regenerated.

With both edges enabled the same callback fires twice per press, so the two
cases must be told apart by reading the current pin level. The button is
active low, so a level of \texttt{GPIO\_PIN\_RESET} means the press has just
occurred and \texttt{GPIO\_PIN\_SET} means the button has just been released.
On press, the current tick is recorded; on release, the elapsed time is
compared against the one-second threshold and the corresponding message is
posted to the queue.

The timestamp must survive between two separate invocations of the interrupt
handler, so it is a file-scope \texttt{volatile} variable rather than a local.

\subsection{Why a Queue Rather Than a Semaphore}

A binary semaphore carries no payload. It can only communicate the fact that
\emph{something} happened, which is sufficient when there is exactly one kind
of event. Once the button can produce two distinguishable events, the
signalling mechanism must carry a value. A message queue provides exactly
this: the interrupt posts a value, and the receiving task both blocks until a
message arrives and learns which event it was. The queue therefore subsumes
the semaphore's role, and Task\_1's binary semaphore was removed entirely.

\subsection{Implementation}

\begin{lstlisting}[caption={EXTI callback: classifying press duration and posting a message.}]
void HAL_GPIO_EXTI_Callback(uint16_t GPIO_Pin)
{
  if (GPIO_Pin == GPIO_PIN_13)
  {
    if (HAL_GPIO_ReadPin(GPIOC, GPIO_PIN_13) == GPIO_PIN_RESET)
    {
      /* falling edge: button pressed, record the time */
      button_press_time = HAL_GetTick();
    }
    else
    {
      /* rising edge: button released, classify the press */
      if (HAL_GetTick() - button_press_time >= 1000)
      {
        uint32_t message = 1;   /* long press  */
        osMessageQueuePut(myQueue01Handle, &message, 0, 0);
      }
      else
      {
        uint32_t message = 0;   /* short press */
        osMessageQueuePut(myQueue01Handle, &message, 0, 0);
      }
    }
  }
}
\end{lstlisting}

The final argument of \texttt{osMessageQueuePut()} is the timeout, and it must
be zero when the function is called from interrupt context. The CMSIS-RTOS v2
wrapper detects that it is running inside an interrupt and rejects any
non-zero timeout outright, because an interrupt handler cannot block.

Task\_1 blocks on the queue and selects its blinking pattern from the received
value:

\begin{lstlisting}[caption={Task\_1: selecting a blink pattern from the queue message.}]
void StartTask03(void *argument)
{
  for(;;)
  {
    uint32_t message;
    osMessageQueueGet(myQueue01Handle, &message, NULL, osWaitForever);

    if (message == 1)          /* long press: 5 s at 10 Hz */
    {
      osMutexAcquire(myMutex01Handle, osWaitForever);
      for (int i = 0; i < 50; i++)
      {
        HAL_GPIO_TogglePin(GPIOB, LED2_Pin);
        osDelay(50);
        HAL_GPIO_TogglePin(GPIOB, LED2_Pin);
        osDelay(50);
      }
      osMutexRelease(myMutex01Handle);
    }
    else                       /* normal press: 5 s at 1 Hz */
    {
      osMutexAcquire(myMutex01Handle, osWaitForever);
      for (int i = 0; i < 5; i++)
      {
        HAL_GPIO_TogglePin(GPIOB, LED2_Pin);
        osDelay(500);
        HAL_GPIO_TogglePin(GPIOB, LED2_Pin);
        osDelay(500);
      }
      osMutexRelease(myMutex01Handle);
    }
  }
}
\end{lstlisting}

\subsection{Results}

A short press produces five slow blinks over five seconds; holding the button
for more than one second and then releasing produces fifty fast blinks over
the same five seconds. In both cases the periodic background burst from
Task\_2 is correctly excluded for the duration of the sequence.

One known limitation is that no debouncing is performed. Mechanical contact
bounce can in principle generate spurious edge interrupts at the moment of
press or release, which would corrupt a duration measurement. In practice the
one-second threshold is long enough that misclassification was not observed,
but a production implementation would filter the input, either in software or
with an RC network.

% =========================================================================
\section{Problems Encountered}
\label{sec:problems}

The majority of the effort in this exercise went into diagnosing behaviour
that appeared to be "nothing happening". Each case had a distinct cause, and
together they illustrate most of the important constraints of programming
under an RTOS.

\subsection{A misspelled callback silently does nothing}

The first version of the interrupt handler was named
\texttt{HAL\_GPIO\_EXIT\_Callback} instead of
\texttt{HAL\_GPIO\_EXTI\_Callback}. The code compiled and linked without any
warning, but the button had no effect.

The HAL declares its callbacks as \emph{weak} symbols with empty bodies, so
that the application can override them simply by defining a function of the
same name. A misspelled name does not override anything; it merely defines an
unused function, and the empty default implementation continues to be called.
The lesson is that a silent failure of this kind is characteristic of the weak
symbol mechanism, and that callback names must be copied exactly.

\subsection{Busy-waiting in a high-priority task starves everything below it}

An early version used \texttt{HAL\_Delay()} inside both tasks. The
higher-priority task blinked correctly while the lower-priority task never ran
at all.

\texttt{HAL\_Delay()} does not block in the RTOS sense: it spins in a loop
polling \texttt{HAL\_GetTick()} and therefore remains permanently
\emph{runnable}. Under a preemptive priority scheduler, a runnable
high-priority task is always selected over a lower-priority one, so the
lower-priority task is never scheduled. Replacing \texttt{HAL\_Delay()} with
\texttt{osDelay()} fixed this: \texttt{osDelay()} moves the calling task into
the blocked state, allowing the scheduler to run other tasks during the delay.

\subsection{\texttt{HAL\_Delay()} inside an interrupt handler deadlocks}

A later attempt performed the blinking loop directly inside the EXTI callback
using \texttt{HAL\_Delay()}. Pressing the button caused the LED to switch on
and freeze permanently.

The cause is an interaction between interrupt priorities. The EXTI interrupt
is configured at NVIC priority 5, whereas the SysTick interrupt that
increments the HAL tick runs at priority 15 --- a lower priority on Cortex-M,
where a numerically larger value is less urgent. While the EXTI handler is
executing, SysTick cannot preempt it, so \texttt{HAL\_GetTick()} never
advances. \texttt{HAL\_Delay()} waits for a tick value that can only be
produced by an interrupt it is itself blocking: a deadlock.

\subsection{\texttt{osDelay()} inside an interrupt handler is a silent no-op}

Replacing \texttt{HAL\_Delay()} with \texttt{osDelay()} in the same handler
removed the freeze but produced no visible blinking at all.

Inspection of the CMSIS-RTOS v2 wrapper shows that \texttt{osDelay()} begins by
testing whether it is executing in interrupt context, and if so returns
\texttt{osErrorISR} immediately without delaying. Every delay in the loop
therefore became a no-op, the ten toggles executed within a few microseconds,
and since ten is an even number the LED finished in the state it started in.
Blocking operations simply cannot be performed in an ISR, regardless of which
API is used.

Taken together, the two previous problems establish the rule that an interrupt
handler must only \emph{signal} a task and return, never perform timed work
itself.

\subsection{A FreeRTOS task function must never return}

While restructuring the code, one task body was temporarily commented out,
including its enclosing \texttt{for(;;)} loop. The whole system stopped
responding.

A task function is not an ordinary function that the scheduler calls
repeatedly; it is an entry point that must never return. When it does,
execution continues into whatever instructions follow it in memory, which in
this case caused a fault and left the CPU spinning in the default
\texttt{HardFault\_Handler}, halting every task. Because the affected task had
the highest priority, it was scheduled first and the system failed almost
immediately after \texttt{osKernelStart()}.

A related misconception is worth recording: calling a task's function by name
from an interrupt handler does not start that task. It executes the function
inline, in interrupt context, on the interrupt's stack --- and if the function
contains an infinite loop, the handler never returns at all.

% =========================================================================
\section{Design Discussion}

\subsection{Signalling and mutual exclusion are different problems}

It is tempting to treat semaphores and mutexes as interchangeable, since both
can be used to make a task wait. They solve different problems. The semaphore
(or queue) answers \emph{``has an event occurred?''} and is signalled by one
party and consumed by another. The mutex answers \emph{``may I touch the
shared resource?''} and is always released by the same task that acquired it.
This design needs both, for different reasons, and the two cannot be merged
without losing one of the guarantees.

An alternative to the mutex would have been a plain \texttt{volatile} flag
checked before each toggle. This works for a single boolean on Cortex-M, where
aligned word accesses are atomic, but it requires the other task to poll the
flag rather than block, and it offers no protection for anything more complex
than a single variable. The mutex expresses the intent directly and lets the
RTOS handle the blocking.

\subsection{Mutex granularity determines takeover latency}
\label{sec:granularity}

A mutex provides mutual exclusion, not preemption. It cannot interrupt a task
that already holds the lock; a waiting task can only be admitted once the
holder releases. The consequence is that the \emph{size of the critical
section} directly determines how quickly one task can take over from another.

In the first implementation, Task\_2 acquired the mutex once and held it for
its entire 2-second burst. A button press arriving early in that burst had to
wait up to two full seconds before Task\_1's sequence could begin, which was
clearly visible as a delayed and apparently unresponsive button.

Narrowing the critical section so that the mutex is held only around each
individual \texttt{HAL\_GPIO\_TogglePin()} call, and released before the
\texttt{osDelay()}, reduced the worst-case wait to approximately one toggle
period (50\,ms) --- indistinguishable from instantaneous to a human observer.
Task\_1 still holds the mutex for its whole sequence, so once it takes over it
cannot be interrupted in turn.

The resulting behaviour is best described as \emph{pause and resume}: Task\_2's
loop counter is preserved while it waits for the mutex, so after Task\_1
finishes, Task\_2 completes whatever remained of its interrupted burst rather
than restarting. Discarding the remainder instead would require an explicit
cancellation mechanism, such as a flag checked inside Task\_2's loop, which was
judged unnecessary for this specification.

\subsection{Sharing SysTick between the HAL and the RTOS}

STM32CubeMX warns that when an RTOS is used, the HAL timebase should be moved
off SysTick onto a separate hardware timer. The reason is that FreeRTOS takes
over the SysTick interrupt for its own scheduling tick, while the HAL relies on
the same interrupt to maintain \texttt{HAL\_GetTick()}. Inside a FreeRTOS
critical section the tick interrupt may be delayed, so HAL timeouts and delays
can measure time inaccurately. This project retains the default shared
configuration, which is acceptable here because no HAL timing function is used
after the scheduler starts, but it would be worth changing in any design that
relies on HAL peripheral timeouts at runtime.

% =========================================================================
\section{Conclusion}

Two tasks were implemented that drive the same LED with different blinking
patterns, one triggered by an external interrupt and one by a periodic RTOS
timer, with a mutex guaranteeing that the two sequences never overlap. The
extension added press-duration measurement using both edges of the EXTI line
and replaced the binary semaphore with a message queue, allowing a single task
to select between two blinking patterns based on the value it receives.

The most valuable outcome was not the final code but the set of failure modes
encountered along the way. Each of them --- the unoverridden weak callback, the
starvation caused by busy-waiting, the deadlock caused by delaying inside a
high-priority interrupt, the silently ignored \texttt{osDelay()} in an ISR, and
the task that returned --- reflects a specific property of how an RTOS
schedules work. Collectively they reinforce a single design principle:
interrupt handlers should capture events and defer all real work to tasks,
which are the only context in which blocking, timing and resource arbitration
can be performed correctly.

\end{document}
